\Chapter{Proposed Electronic Medical Record system solutions}\label{chapter:project-idea}

\setlength{\parskip}{0pt}
\setlength{\parindent}{1.25cm}

\section{Introduction to the design process}
Based on the analysis of the operation of Szpital im. Bajgla, several possible directions for solving the problems related to the organization of medical documentation were prepared. Due to the scale of the facility and the complexity of medical processes, there is no single universal solution that would respond in the same way to all organizational, technical and financial needs of the hospital.

The design process was therefore based on the analysis of various EMR (Electronic Medical Record) implementation options, differing in the scope of functionality and the level of technological advancement. The prepared solutions take into account both the current needs of the hospital and the possibility of further system development in the future.

An important aspect of the design process was also taking into account the needs of various user groups, including doctors, nurses, laboratories, administration and patients themselves.

The prepared variants cover three levels of system implementation: a basic solution focusing on the digitalization of the most important processes, an optimal solution constituting a complete integrated EMR system for a large hospital, and an extended solution covering additional automation and intelligent facility management functions.


\section{Design assumptions}
Each of the analyzed system variants was developed on the basis of design assumptions and requirements related to the operation of modern healthcare systems. Regardless of the implementation scope, all solutions must ensure the possibility of central storage of medical documentation and handling multiple departments within a single IT environment.

An important design assumption is also the adaptation of the system to the work of various user groups. The interface should be intuitive and as simple as possible in everyday use, so as to reduce the risk of errors resulting from improper use. The project also assumes the possibility of using the system on desktop or mobile devices used by medical staff while working in departments.

Each of the variants must also meet the requirements concerning medical data security and compliance with regulations in force in the European Union. The system should enable control of access to documentation, separation of user permissions, monitoring of activity in the system and protection of data against loss or unauthorized access.
\section{Variant I - basic solution}

\subsection{Characteristics of the solution}
The basic variant assumes the digitalization of the most important processes related to maintaining medical documentation at Szpital im. Bajgla. The solution focuses primarily on moving away from paper-based storage of patient data and creating a central place enabling the viewing of basic medical documentation by hospital staff.
The system in this version is simplified in nature and includes only the most important functions necessary to improve the organization of the facility's work.

The project assumes the use of computer workstations located in departments without implementing an extensive mobile infrastructure. This variant was prepared as a solution requiring relatively small organizational changes and a limited implementation scope.

The basic solution focuses mainly on reducing the problems associated with storing paper documentation and simplifying access to patient information. However, the system does not include advanced process automation or full integration of all elements of the hospital's operation.
\subsection{Scope of functionality}
The basic variant includes a central patient database enabling the storage of basic medical documentation in electronic form. Medical staff can view the patient's treatment history, test results and basic information concerning hospitalization from computer workstations located in departments.

The system also enables the electronic addition of laboratory test results and diagnostic documentation. The solution provides a basic user login model and a division of permissions between medical staff and administration.

The variant also includes a simplified patient portal enabling access to selected test results and basic information concerning treatment.

\subsection{Advantages of the solution}
The greatest advantage of the basic variant is the relatively low level of implementation complexity and limited infrastructure requirements. The solution allows a quick reduction of some of the problems resulting from paper documentation without the need for a complete reorganization of the hospital's work.

The basic variant may also reduce the resistance of some staff to digitalization, because the scope of organizational changes remains relatively small. The system is simpler to maintain and requires fewer training sessions than more extensive solutions.

An additional advantage is the possibility of gradually developing the system in the future by adding further functions along with the facility's further digitalization.

\subsection{Limitations of the solution}
The basic variant does not completely eliminate the problems related to information flow between departments and the limited mobility of medical staff. The lack of tablet use may cause further delays in access to data in situations requiring a rapid response.

Another limitation of the solution is the low level of automation of administrative processes and the lack of advanced functions supporting the organization of work of a large medical facility. The system does not ensure full integration of all hospital processes, which means that some activities will still require manual transfer of information between users.

The basic variant may also prove to be an insufficient solution in the perspective of further development of the facility and the growing requirements related to the digitalization of healthcare.

\section{Variant II - optimal solution}

\subsection{Characteristics of the solution}
The optimal variant assumes the implementation of a complete integrated EMR system covering all key areas of the hospital's operation. The solution was designed for a large multi-department facility requiring efficient information flow between medical staff, laboratories, diagnostics and administration.

The system in this version is to constitute a central working environment for the entire hospital, enabling ongoing access to patient documentation in real time. The project assumes the use of both computer workstations and mobile devices used by medical staff while working in departments.

The solution includes full digitalization of medical documentation and automation of selected organizational processes related to the circulation of information and test results.

\subsection{Scope of functionality}
The optimal variant includes a central patient database, electronic medical records, electronic test orders and ongoing data synchronization between departments. Medical staff can view and update patient data in real time from computers and tablets.

The system enables automatic transfer of laboratory and diagnostic test results to the patient's documentation and monitoring of the status of performed tests. The solution also provides a notification module informing staff about new results, documentation changes or important events related to patient treatment.

The project includes an extensive system of roles and access levels, documentation change history, data backup and mechanisms for monitoring user activity. Patients have access to a portal enabling them to check test results, treatment history and selected information concerning hospitalization.


\subsection{Advantages of the solution}
The greatest advantage of the optimal variant is the full integration of processes related to medical documentation and information flow between departments. The system significantly improves communication between medical staff.

The solution improves the organization of the hospital's work by enabling rapid access to current patient data. The use of mobile devices increases the convenience of staff work and enables access to documentation without the need to return to stationary computer workstations.

The optimal variant also ensures a high level of medical data security and the possibility of further system expansion in the future. The solution meets the needs of a large medical facility and remains consistent with the direction of development of modern eHealth systems in the European Union.

\subsection{Limitations of the solution}
Implementing the optimal variant requires greater organizational and technical effort than the basic solution. It becomes necessary to prepare appropriate network infrastructure, purchase mobile devices and conduct training for medical staff.

The solution also involves a greater level of organizational changes in the facility's daily functioning. Some employees may initially experience difficulties related to the transition to a fully digital working model.

The variant also requires more extensive technical support and regular system maintenance due to the greater number of functionalities and users.

\section{Variant III - extended solution}

\subsection{Characteristics of the solution}
The extended variant assumes the expansion of the EMR system with additional data analysis functions and partial automation of the organization of the hospital's work. Unlike the optimal variant, the solution does not focus solely on electronic documentation and communication between departments, but additionally enables ongoing monitoring of the facility's operation and analysis of organizational data.

The system is to support the management and administration in controlling the functioning of the hospital through the use of more advanced reporting, monitoring and analysis tools for data generated during the facility's daily work.
\subsection{Scope of functionality}
The extended variant includes additional analytical modules and functions for automating organizational processes. The system can automatically create reports concerning department workload, the number of patients, waiting times for tests or the use of medical staff.

The solution can also generate intelligent notifications concerning organizational problems, such as test delays, department overload or missing documentation elements. In addition, the possibility of integration with selected medical devices was provided, thanks to which some data could be transferred to the system automatically without the need for manual rewriting of results.

The project also includes a more extensive patient application enabling broader contact with the hospital and access to a larger amount of information concerning treatment.
\subsection{Advantages of the solution}
The greatest advantage of the extended variant is the possibility of using data to support the organization of the hospital's work and faster detection of problems related to the facility's operation. The system can facilitate the analysis of department workload and improve the management of staff and hospital resources.

An additional advantage is the possibility of further expanding the system with subsequent automation functions and integration with additional devices and modules in the future.
\subsection{Limitations of the solution}
The extended variant requires more extensive technical infrastructure and greater administrative support than the optimal variant. The solution also involves a greater level of complexity in maintaining the system and the need for more advanced management of its operation.

In addition, in the case of integrating data analysis mechanisms based on artificial intelligence, it would be necessary to adapt the system to the requirements resulting from the AI Act and regulations concerning personal data protection. This would mean, among other things, the need to properly document the system's operation, monitor the use of AI models and ensure transparency of the medical data processing process.

It would also be necessary to inform patients and employees about the use of artificial intelligence mechanisms in the analysis of medical and organizational data. Depending on the scope of the system's operation, it might also be required to implement additional procedures related to supervision over the operation of AI models and data security.

\section{Comparison of variants}

\subsection{Comparison table}

\renewcommand{\arraystretch}{1.3}

\begin{longtable}{|p{4cm}|p{3cm}|p{3cm}|p{3cm}|}

\caption{Comparison of EMR system variants}
\label{tab:warianty_emr} \\

\hline
\textbf{Criterion} &
\shortstack{\textbf{Variant I} \\ basic} &
\shortstack{\textbf{Variant II} \\ optimal} &
\shortstack{\textbf{Variant III} \\ extended} \\
\hline
\endfirsthead

\hline
\textbf{Criterion} &
\shortstack{\textbf{Variant I} \\ basic} &
\shortstack{\textbf{Variant II} \\ optimal} &
\shortstack{\textbf{Variant III} \\ extended} \\
\hline
\endhead

Scope of documentation digitalization &
partial &
full &
full \\
\hline

Central patient database &
yes &
yes &
yes \\
\hline

Electronic medical record &
basic &
full &
full \\
\hline

Department integration &
limited &
full &
full \\
\hline

Real-time data access &
partial &
yes &
yes \\
\hline

Mobile access for staff &
none &
tablets and mobile devices &
tablets and mobile devices \\
\hline

Patient portal &
basic &
extended &
extended \\
\hline

Automatic transfer of test results &
partial &
yes &
yes \\
\hline

Analytics modules and reporting &
none &
basic reports &
advanced data analysis \\
\hline

Department workload analysis &
none &
none &
extensive \\
\hline

Smart notifications and alerts &
none &
basic &
extensive \\
\hline

Integration with medical devices &
none &
limited &
extended \\
\hline

Support for hospital work organization &
basic &
high &
very high \\
\hline

Data security and access control &
basic &
high &
very high \\
\hline

Use of AI data analysis &
none &
none &
yes \\
\hline

AI Act requirements &
not applicable &
not applicable &
required \\
\hline

Level of automation &
low &
medium &
high \\
\hline

Implementation complexity &
low &
high &
very high \\
\hline

Infrastructure requirements &
low &
medium &
high \\
\hline

Scope of staff training &
small &
large &
very large \\
\hline

\end{longtable}

\subsection{Comparative analysis}

The conducted analysis shows that each of the variants responds to the hospital's needs to a different extent. The basic variant enables partial digitalization of documentation and a reduction in the amount of paper documents, but it does not ensure full integration of medical and administrative processes. This solution is characterized by the lowest implementation costs and the lowest level of infrastructure complexity.

The optimal variant ensures full digitalization of medical documentation, integration of departments and efficient access to patient data in real time. The solution offers a wide range of functionalities while maintaining moderate organizational and infrastructure requirements.

The extended variant offers the highest level of automation and additional functions related to data analysis and the integration of artificial intelligence elements. However, the solution involves the greatest level of implementation complexity, high maintenance costs and additional legal requirements related, among other things, to the AI Act.

\section{Selection of the target solution}

Based on the conducted analysis, the optimal variant was selected as the target solution. This solution constitutes a compromise between the scope of functionality, implementation costs and the organizational capabilities of the hospital.

The selected variant meets all the key requirements of the facility related to full digitalization of medical documentation, improvement of information flow and increased availability of patient data. At the same time, it does not require such extensive infrastructure and such large organizational changes as the extended variant.

The client decided on a solution that is feasible to implement under the hospital's current financial conditions, while ensuring the possibility of further system expansion in the future.
\section{Solution evaluation criteria}
After the introduction of the target EMR system variant, an assessment of the effectiveness of the implemented solution will be carried out based on defined functional, organizational and technical criteria. The purpose of the assessment will be to check to what extent the system improves the functioning of the hospital and reduces the problems associated with paper-based documentation circulation.

The first criterion will be the level of digitalization of medical documentation. It will be determined as the percentage of documents handled electronically in relation to all documentation created in the hospital's daily work. This will include, among others, patient charts, test results, medical orders, nursing reports, discharge summaries and administrative documents. The greater the share of documents handled in the digital system, the higher the level of digitalization of the facility.

The second criterion will be the time of access to patient documentation. The measurement will be carried out by checking how much time an authorized user needs to find the treatment history, current test results or recent medical orders. The analysis will cover typical hospital work situations, such as patient admission, medical rounds or preparation of a discharge summary.

The third criterion will be the availability of the system for users in daily work. The assessment will cover the number of workstations and devices enabling the use of the system, the possibility of mobile work and the availability of data in departments where staff do not work permanently at a computer. It will be particularly important to check whether a doctor or nurse can access patient data without having to leave the department or interrupt the activities being performed.

The fourth criterion will be security and control of access to data. The assessment will check whether the system has user roles, individual accounts, documentation change history, a login register and the ability to limit access to data in accordance with the employee's scope of duties. This criterion will be measured, among other things, by the number of available permission levels and the ability to monitor user activity in the system.
